\scanpage{163}
\mspage{488}
% position: upper continuation paragraph
que j'avais parlé à contre temps, faute d'écoute et de discernement, avant que
cette propension invétérée commence à me passer. Si elle est devenue beaucoup
moins envahissante, et qu'elle a même quasiment disparue, ce n'est nullement \unsure{com|com-}\linebreak
me résultat de quelque discipline que je me serais imposé\inserted{\unsure{e}} (style: tu n'ouvriras
le bec que si...). C'est simplement parce que l'envie m'a passé de parler, en \unsure{de|des}
moments où je sens bien que c'est inutile, que cela n'apporte rien à autrui ni
à moi - rien du moins qui ait une valeur à mes yeux. Si à présent j'arrive sou-\linebreak
vent à sentir de telles choses, c'est sans doute que je suis devenu plus atten-\linebreak
tif. Cela aussi n'est pas venu comme fruit d'une discipline (``tu vas faire
gaffe d'ouvrir grandes tes oreilles quand...''), mais je ne saurais dire comment.
En tous cas, je m'en sens mieux, la vie est d'autant plus intéressante (et sur-\linebreak
tout moins bruyante!). Et les autres aussi s'en sentent mieux...

% position: lower paragraph
Je crois que j'ai commencé à vraiment moins parler, à partir du moment
où à disparu (autant dire) cette force en moi qui me pousse à toujours vouloir
rectifier ce qui m'apparaît (à tort ou à raison) comme des ``erreurs'' chez au-\linebreak
trui - comme s'il n'était pas suffisant que je décèle et rectifie les miennes!
C'est la force aussi qui me poussait (et parfois encore me pousse) à vouloir
à toute fin convaincre autrui de ceci ou de celà, au lieu de regarder simplement
pourquoi untel préfère mordicus croire ceci plutôt que cela (qui me paraît ``ça'',
et dont je voudrais bien le convaincre!); ou pourquoi je tiens tellement à ce
\unsure{qu'il} croie cela, plutôt que ceci. Cette force quasiment universelle en nous,
qui nous pousse constamment à chercher dans l'approbation d'autrui (et ne fut-ce
qu'un seul...) la confirmation du bien-fondé de ce que nous tenons pour vrai -
cette force profondément enracinée dans l'égo a fini, je crois, par lâcher pri-\linebreak
se en moi. Cela a été un grand soulagement, la fin d'une dispersion d'énergie
faramineuse. C'est quand j'ai fini par me rendre compte, il y a deux ans, de la
portée de cette force dans ma vie, de sa nature, et de d'extraordinaire disper-\linebreak
sion d'énergie qu'elle représentait, qu'elle s'est trouvée désamorcée - et que
je me suis trouvé allégé du coup ``d'un poids de cent tonnes''. Prendre connais-\linebreak
sance sans réticence de d'écho qu'autrui nous renvoie de notre personne, sans
être lié par un désir ou ``besoin'' (si caché soit-il) d'approbation ou de confir-\linebreak
mation - c'est cela vraiment, être ``libre de lui''. C'est un tel besoin ou désir
qui constitue véritablement le ``crochet'', discret et d'une solidité à toute
épreuve, par où le conflit peut ``accrocher'' en nous, et par où nous sommes (que
nous le voulions ou le reconnaissions, ou non) sous la dépendance d'autrui, de
% position: left lower mark
\margin{\unsure{-}}
% position: foot pencil
\margin{162}
